% Answer sheet for the Chapter 2 exercises. Build here: pdflatex (run twice).
\documentclass[12pt]{article}
\newcommand{\guideroot}{../}
% Shared preamble for the BARCS student guide (both chapters and the combined book).
% Written to the rules of ASD-STE100 Simplified Technical English:
% short sentences, active voice, present tense, one idea per sentence,
% one meaning per word, and every symbol defined before use.
\usepackage{needspace}
\usepackage[margin=1in]{geometry}
\usepackage[T1]{fontenc}
\usepackage{lmodern}
\usepackage{amsmath,amssymb}
\usepackage{booktabs}
\usepackage{array}
\usepackage{enumitem}
\usepackage[most]{tcolorbox}
\usepackage{xcolor}
\usepackage{tikz}
\usetikzlibrary{arrows.meta,positioning}
\usepackage[hidelinks]{hyperref}

\definecolor{barcsblue}{HTML}{0072B2}
\definecolor{barcsgreen}{HTML}{009E73}
\definecolor{barcsdark}{HTML}{1F3B4D}
\definecolor{barcsorange}{HTML}{D55E00}

\newtcolorbox{symbols}{colback=barcsblue!5,colframe=barcsblue,
  fonttitle=\bfseries,title=What the symbols mean,boxrule=0.6pt,arc=2pt}
\newtcolorbox{worked}[1][]{colback=barcsgreen!5,colframe=barcsgreen!80!black,
  fonttitle=\bfseries,title=Worked example#1,boxrule=0.6pt,arc=2pt}
\newtcolorbox{remember}{colback=barcsorange!5,colframe=barcsorange,
  fonttitle=\bfseries,title=Remember,boxrule=0.6pt,arc=2pt}

\newlist{choices}{enumerate}{1}
\setlist[choices]{label=(\Alph*),leftmargin=2.2em,itemsep=1pt,topsep=2pt}
\setlist{itemsep=2pt}

\setlength{\parskip}{6pt}
\setlength{\parindent}{0pt}
\setcounter{secnumdepth}{0}  % the steps carry their own numbers
\setcounter{tocdepth}{2}
\newif\ifstandalone
\newcommand{\guideauthor}{Prepared for Hyun-Hwan Jeong and Kyu-Won Lee}
\newcommand{\guidedate}{BARCS version 0.2.1 \quad$\cdot$\quad Reading level: grades 7 to 9}
\newcommand{\aboutbox}{%
\begin{center}
\begin{minipage}{0.85\textwidth}
\small
\textbf{About this guide.} This guide explains how the BARCS computer
program works. You need only arithmetic, fractions, percents, powers and
square roots. The guide teaches logarithms when it uses them. Both chapters
use the same small set of numbers, so you can follow every step. The guide
uses the writing rules of ASD-STE100 Simplified Technical English.
\end{minipage}
\end{center}}
% Divider page that opens a chapter in the combined book.
\newcommand{\chapterpart}[3]{%
  \clearpage
  \addcontentsline{toc}{part}{Chapter #1: #2}%
  \vspace*{0.25\textheight}
  \begin{center}
  {\Huge\bfseries Chapter #1\par}\vspace{8pt}
  {\LARGE #2\par}\vspace{6pt}
  {\large #3\par}
  \end{center}
  \clearpage}
\providecommand{\guideroot}{}
% Exercise tags: what kind of thinking a question asks for.
\newcommand{\tagcalc}{\textcolor{barcsblue}{\textsf{\footnotesize[Calculate]}}}
\newcommand{\tagexplain}{\textcolor{barcsgreen!70!black}{\textsf{\footnotesize[Explain]}}}
\newcommand{\tagdecide}{\textcolor{barcsorange}{\textsf{\footnotesize[Decide]}}}
\newcommand{\tagwrite}{\textcolor{barcsdark}{\textsf{\footnotesize[Write]}}}
\newlist{exlist}{enumerate}{1}
\setlist[exlist]{label=\textbf{\arabic*.},leftmargin=2em,itemsep=8pt}
\newlist{parts}{enumerate}{1}
\setlist[parts]{label=(\alph*),leftmargin=2em,itemsep=2pt,topsep=2pt}

\standalonetrue
\title{\textbf{How BARCS Finds Important Genes}\\[6pt]
\Large Answer sheet: Chapter 2 exercises}
\author{\guideauthor}
\date{\guidedate}
\newcommand{\why}[1]{\par\smallskip\noindent\textcolor{barcsgreen!70!black}{\textbf{Why this matters.}} #1}
\begin{document}
\maketitle
\thispagestyle{empty}

Each answer shows the calculation first. ``Why this matters'' explains the
idea behind the question. Values: $\mu = 0.00015$ (control), $\mu =
0.000075$ (drug), $\rho = 0.000001$, $\beta_1 = -0.69$.

\section*{Part A: Fitting}
\begin{exlist}
\item
\begin{parts}
\item $0.00015 \times 0.99985 = 0.000149978$. Slope: $1 \div 0.000149978 =
  6{,}668$.
\item Step: $(0.00018 - 0.00015) \times 6{,}668 = 0.00003 \times 6{,}668 =
  0.200$. Working value: $v = -8.8047 + 0.200 = -8.6047$.
\item $0.00018 \div 0.99982 = 0.000180032$. $\ln(0.000180032) = -8.6224$.
  The difference is only $0.018$.
\item With $K = 0$, the real log-odds is $\ln(0)$, which does not exist. The
  working value still exists: $v = -8.8047 + (0 - 0.00015) \times 6{,}668 =
  -8.8047 - 1.000 = -9.805$. The straight line also makes the variance of
  $v$ easy to calculate, which gives the weight.
\end{parts}
\why{The working value lets BARCS use simple weighted averages, even for
guides with zero reads.}

\item
\begin{parts}
\item Library with 1 million reads: top $= 1{,}000{,}000 \times 0.00015
  \times 0.99985 = 149.98$; bottom $= 1 + 999{,}999 \times 0.000001 = 2.0$;
  $w = 149.98 \div 2.0 = 75.0$.\\
  Library with 2 million reads: top $= 299.96$; bottom $= 1 + 1{,}999{,}999
  \times 0.000001 = 3.0$; $w = 299.96 \div 3.0 = 100.0$.
\item No. Twice the reads gives only $100 \div 75 = 1.33$ times the weight.
  The extra reads shrink the counting noise, but the biological noise of
  the flask stays. The bottom of the weight grows from 2 to 3.
\item With $\rho = 0$, the bottom is 1. The weights are 150 and 300, exactly
  double.
\end{parts}
\why{The weight says how much a library can tell you. Depth helps only
until biological noise dominates.}

\item
\begin{parts}
\item $(75 \times (-8.60) + 100 \times (-8.90)) \div (75 + 100) = (-645 -
  890) \div 175 = -1{,}535 \div 175 = -8.771$.
\item $\sqrt{1 \div 175} = \sqrt{0.005714} = 0.0756$.
\item The weight is $1 \div$ variance. The library with weight 100 has less
  noise, so it should count more. ``Weight $= 1 \div$ variance'' gives the
  average with the smallest possible variance.
\end{parts}
\why{A weighted average trusts the precise libraries more. This is the core
of the BARCS fit.}
\end{exlist}

\section*{Part B: Is the effect real?}
\begin{exlist}[resume]
\item
\begin{parts}
\item $1 \div 175 = 0.005714$; $1 \div 70 = 0.014286$; sum $= 0.020000$;
  $\text{SE} = \sqrt{0.020000} = 0.1414$.
\item $t = -0.69 \div 0.1414 = -4.88$. $\text{df} = 4 - 2 = 2$.
\item $|t| = 4.88$ is larger than 4.30. So $p < 0.05$ (a computer gives
  $p = 0.040$). It is significant, but only just.
\item The drug group: $0.014286$ of the $0.020000$ is from the drug group
  (71\%). Its total weight is small because its libraries are smaller and
  its proportion is lower. Add drug replicates. That adds weight where it is
  missing, and it adds degrees of freedom.
\end{parts}
\why{The SE shows where the uncertainty comes from. That tells you how to
design the next experiment.}

\item
\begin{parts}
\item Guide P: $5.1 \div 2 = 2.55$. Guide Q: $1.4 \div 2 = 0.70$.
\item Guide P has a sum above 2, so it has more noise than counting noise:
  BARCS uses $\rho > 0$. Guide Q has a sum of 2 or less: BARCS uses
  $\rho = 0$.
\item The predictions use two $\beta$ values that were fitted to the same
  four libraries. The fit pulls the predictions toward the data, and each
  fitted $\beta$ takes away 1 from the sum on average. So the expected sum
  is $4 - 2 = 2$, the degrees of freedom.
\end{parts}
\why{BARCS measures $\rho$ by making the noise model agree with the data on
average.}

\item
\begin{parts}
\item $(2 \times 2.55 + 10 \times 1.5) \div (2 + 10) = (5.1 + 15) \div 12 =
  20.1 \div 12 = 1.675$.
\item $\sqrt{1.675 \div 2.55} = \sqrt{0.657} = 0.81$. The SE becomes 19\%
  smaller. New degrees of freedom: $2 + 10 = 12$. Now $|t|$ needs to be
  larger than only 2.18 (not 4.30) for $p < 0.05$.
\item The guides of one screen went through the same cells, the same
  library preparation and the same sequencer. So their noise is similar. One
  guide with 2 degrees of freedom gives a very unsure noise estimate, and
  the other guides make it steadier. It can be a bad idea when a guide
  really is much noisier than the rest, for example when a few cell clones
  take over a flask. Then borrowing makes that guide look too precise. The
  check with control guides (Exercise~7) catches this problem.
\end{parts}
\why{Moderation turns thousands of guides into extra information about
noise. That is how a small experiment can still find real effects.}
\end{exlist}

\section*{Part C: From guides to a list of genes}
\begin{exlist}[resume]
\item
\begin{parts}
\item $0.05 \times 300 = 15$. The real number is 45, three times more. The
  noise model is too small: the libraries are noisier than the model
  thinks. One common cause is replicates that are not independent, for
  example flasks split from one population of cells.
\item $2.6 \div 1.96 = 1.33$.
\item $3.0 \div 1.33 = 2.26$. It is still beyond 1.96, so the gene keeps
  $p < 0.05$ (with many degrees of freedom).
\end{parts}
\why{Control guides have no true effect. They are a free test of the noise
model inside your own experiment.}

\item
\begin{parts}
\item $(2.2 + 1.8 - 0.3 + 2.5) \div \sqrt{4} = 6.2 \div 2 = 3.1$
  ($p = 0.002$).
\item Bottom: $\sqrt{4 + 4 \times 3 \times 0.2} = \sqrt{4 + 2.4} =
  \sqrt{6.4} = 2.53$. $Z = 6.2 \div 2.53 = 2.45$ ($p = 0.014$).
\item A negative $z$ means that this guide moved the other way. It is
  evidence against a consistent effect, so it lowers the sum. With shared
  noise, the guides agree partly because they share noise, not because of
  the gene. Each of the $4 \times 3 = 12$ pairs adds $r$ to the variance of
  the sum (squares and pairs, Chapter~1). A bigger variance means a bigger
  bottom and a smaller $Z$.
\end{parts}
\why{A gene is believable when its independent guides agree. Shared noise
makes agreement less convincing.}

\item
\begin{parts}
\item $10{,}000 \times 0.0005 = 5$ genes.
\item About $5 \div 40 = 0.125$. About one gene in eight on this list is
  probably false.
\item Use a smaller limit. Fewer genes pass by chance, and the list usually
  becomes cleaner, but also shorter. Benjamini--Hochberg finds the limit
  for you.
\end{parts}
\why{With thousands of tests, a few small p-values happen by chance. The
FDR tells you how many of your hits to expect to be false.}

\item The extra factor at 20 million reads is $1 + 19{,}999{,}999 \times
0.000001 \approx 21$. So the biological noise is about 20 times the counting
noise. The standard deviation of the proportion is 12.5 per million,
almost the same as the limit of 12.2 per million (Chapter~1, Exercise~8),
and not much better than 17.3 per million at 1 million reads. The deep
reading buys almost nothing. Also, 2 replicates per group give 4 libraries
and only $4 - 2 = 2$ degrees of freedom, so $|t|$ must be larger than 4.30.
A better plan: more replicates per group (for example 4) with fewer reads
each.
\why{Reads and replicates are not the same thing. Degrees of freedom and
biological noise come from replicates only.}

\item
\begin{parts}
\item $(2.0 + 2.6 + 2.2) \div 3 = 6.8 \div 3 = 2.27$.
\item $p = 2 \times (1 + 9) \div (1 + 1{,}999) = 2 \times 10 \div 2{,}000 =
  0.01$.
\item $\sqrt{1 + (3 - 1) \times 0.1} = \sqrt{1.2} = 1.10$. The guides of a
  real gene share some noise, so their average wobbles more than the
  average of unrelated control guides. Without the widening, the
  pseudo-genes would look calmer than a real gene with no effect, and real
  genes would look too significant.
\item $2 \div (150 + 1) = 2 \div 151 = 0.013$. No one-guide gene can get a
  smaller p-value, and the FDR step needs many genes at that floor. So
  ``no hits'' here means ``150 controls are too few to confirm the hits'',
  not ``nothing changed''. Add more control guides to the next library, or
  look at genes with several guides.
\end{parts}
\why{The empirical null trusts the controls of your own experiment. Its
strength (no model needed) is also its limit (it can only be as fine as the
number of controls).}
\end{exlist}
\end{document}
